% GENERATED FILE. Edit canonical lessons, then run npm run generate:books.
% canonical-source-hash: fnv1a64:bd066818d6b9b8a3
% canonical-lessons: HI-C397-visay, HI-C397-patrika, HI-C397-chatravrtti, HI-C397-farq, HI-C397-praptkarna, HI-R397-review-391-397, HI-R397-review-final-397

\chapter{Choosing What to Study}
\label{ch:hi-choosing-what-to-study}
\hlchaptermodality{397}

\begin{chapteropening}
\textbf{By the end of this chapter:} \emph{I can say what I will study if I get the chance.}

Complete the last lesson of chapter 397: I can say what I will study if I get the chance.
\end{chapteropening}

\section[\texorpdfstring{\hi{विषय} (viṣay)}{विषय (viṣay)}]{\hi{विषय} (viṣay) — a subject, a topic}

\label{lesson:HI-C397-visay}



\begin{quote}
Before the new one: say the Hindi for a good deed, then the Hindi for to learn by rote.
\end{quote}

\subsection*{You'll want to know: \hi{विषय}}
\textbf{\hi{विषय}} — \emph{viṣay} — \textquotedblleft{}a subject, a topic\textquotedblright{}.

\begin{quote}
\hi{अगर} \hi{तुम्हें} \hi{विज्ञान} \hi{पसंद} \hi{है} \hi{तो} \hi{वही} \hi{विषय} \hi{चुनो।}
\end{quote}

\begin{quote}
\emph{agar tumhẽ vijñān pasand hai to vahī viṣay cuno.}
\end{quote}

\begin{quote}
If you like science, choose that subject.
\end{quote}

Masculine.

\begin{cousinweb}
\textbf{\hi{विषय}} (\emph{viṣay}) is a tatsama, borrowed from Sanskrit \emph{viṣaya}, \textquotedblleft{}a sphere, a scope\textquotedblright{} (Monier-Williams).
\end{cousinweb}

\begin{grammarlens}[title={What you've built}]
One more word for saying what happens if.
\end{grammarlens}

\subsection*{Your turn}
\begin{itemize}
\raggedright
  \item \emph{Say it:} \emph{viṣay}
  \item \emph{Say it:} \emph{viṣay}, once more
  \item \emph{Say it:} \emph{agar tumhẽ vijñān pasand hai to vahī viṣay cuno.} — If you like science, choose that subject.
  \item \emph{Recall:} say the Hindi for a good deed, then the Hindi for to learn by rote, then say \emph{viṣay} again
\end{itemize}

\begin{tcolorbox}[breakable,colback=teal!4,colframe=teal!35!black,title={Before you move on}]
What does \hi{विषय} mean? (\textquotedblleft{}A subject, a topic\textquotedblright{}.) Say it once more.
\end{tcolorbox}
\section[\texorpdfstring{\hi{पत्रिका} (patrikā)}{पत्रिका (patrikā)}]{\hi{पत्रिका} (patrikā) — a magazine}

\label{lesson:HI-C397-patrika}



\begin{quote}
Before the new one: say the Hindi for to learn by rote, then the Hindi for a subject.
\end{quote}

\subsection*{You'll want to know: \hi{पत्रिका}}
\textbf{\hi{पत्रिका}} — \emph{patrikā} — \textquotedblleft{}a magazine\textquotedblright{}.

\begin{quote}
\hi{अगर} \hi{तुम} \hi{पढ़ोगे} \hi{तो} \hi{यह} \hi{पत्रिका} \hi{तुम्हें} \hi{बहुत} \hi{अच्छी} \hi{लगेगी।}
\end{quote}

\begin{quote}
\emph{agar tum paṛhoge to yah patrikā tumhẽ bahut acchī lagegī.}
\end{quote}

\begin{quote}
If you read it, you'll like this magazine a lot.
\end{quote}

Feminine.

\begin{cousinweb}
\textbf{\hi{पत्रिका}} (\emph{patrikā}) is a tatsama, borrowed from Sanskrit \emph{pattrikā}, \textquotedblleft{}a leaf; a letter, a document\textquotedblright{} (Monier-Williams).
\end{cousinweb}

\begin{grammarlens}[title={What you've built}]
One more word for saying what happens if.
\end{grammarlens}

\subsection*{Your turn}
\begin{itemize}
\raggedright
  \item \emph{Say it:} \emph{patrikā}
  \item \emph{Say it:} \emph{patrikā}, once more
  \item \emph{Say it:} \emph{agar tum paṛhoge to yah patrikā tumhẽ bahut acchī lagegī.} — If you read it, you'll like this magazine a lot.
  \item \emph{Recall:} say the Hindi for to learn by rote, then the Hindi for a subject, then say \emph{patrikā} again
\end{itemize}

\begin{tcolorbox}[breakable,colback=teal!4,colframe=teal!35!black,title={Before you move on}]
What does \hi{पत्रिका} mean? (\textquotedblleft{}A magazine\textquotedblright{}.) Say it once more.
\end{tcolorbox}
\section[\texorpdfstring{\hi{छात्रवृत्ति} (chātravṛtti)}{छात्रवृत्ति (chātravṛtti)}]{\hi{छात्रवृत्ति} (chātravṛtti) — a scholarship (a grant for study)}

\label{lesson:HI-C397-chatravrtti}



\begin{quote}
Before the new one: say the Hindi for a subject, then the Hindi for a magazine.
\end{quote}

\subsection*{You'll want to know: \hi{छात्रवृत्ति}}
\textbf{\hi{छात्रवृत्ति}} — \emph{chātravṛtti} — \textquotedblleft{}a scholarship (a grant for study)\textquotedblright{}.

\begin{quote}
\hi{अगर} \hi{मुझे} \hi{छात्रवृत्ति} \hi{मिलेगी} \hi{तो} \hi{मैं} \hi{शहर} \hi{में} \hi{पढ़ूँगा।}
\end{quote}

\begin{quote}
\emph{agar mujhe chātravṛtti milegī to maiṁ śahar meṁ paṛhūṅgā.}
\end{quote}

\begin{quote}
If I get a scholarship, I'll study in the city.
\end{quote}

Feminine.

\begin{cousinweb}
\textbf{\hi{छात्रवृत्ति}} (\emph{chātravṛtti}) joins two Sanskrit words, \emph{chāttra}, \textquotedblleft{}a pupil\textquotedblright{} (Hindi \hi{छात्र}), and \emph{vṛtti}, \textquotedblleft{}conduct; a livelihood\textquotedblright{} (Monier-Williams): a pupil's keep. The compound is a modern one.
\end{cousinweb}

\begin{grammarlens}[title={What you've built}]
One more word for saying what happens if.
\end{grammarlens}

\subsection*{Your turn}
\begin{itemize}
\raggedright
  \item \emph{Say it:} \emph{chātravṛtti}
  \item \emph{Say it:} \emph{chātravṛtti}, once more
  \item \emph{Say it:} \emph{agar mujhe chātravṛtti milegī to maiṁ śahar meṁ paṛhūṅgā.} — If I get a scholarship, I'll study in the city.
  \item \emph{Recall:} say the Hindi for a subject, then the Hindi for a magazine, then say \emph{chātravṛtti} again
\end{itemize}

\begin{tcolorbox}[breakable,colback=teal!4,colframe=teal!35!black,title={Before you move on}]
What does \hi{छात्रवृत्ति} mean? (\textquotedblleft{}A scholarship (a grant for study)\textquotedblright{}.) Say it once more.
\end{tcolorbox}
\section[\texorpdfstring{\hi{फ़र्क़} (farq)}{फ़र्क़ (farq)}]{\hi{फ़र्क़} (farq) — a difference}

\label{lesson:HI-C397-farq}



\begin{quote}
Before the new one: say the Hindi for a magazine, then the Hindi for a scholarship.
\end{quote}

\subsection*{You'll want to know: \hi{फ़र्क़}}
\textbf{\hi{फ़र्क़}} — \emph{farq} — \textquotedblleft{}a difference\textquotedblright{}.

\begin{quote}
\hi{अगर} \hi{तुम} \hi{रोज़} \hi{एक} \hi{घंटा} \hi{पढ़ोगे} \hi{तो} \hi{बड़ा} \hi{फ़र्क़} \hi{पड़ेगा।}
\end{quote}

\begin{quote}
\emph{agar tum roz ek ghaṇṭā paṛhoge to baṛā farq paṛegā.}
\end{quote}

\begin{quote}
If you study for an hour every day, it will make a big difference.
\end{quote}

Masculine.

\begin{cousinweb}
\textbf{\hi{फ़र्क़}} (\emph{farq}) is usually traced to Arabic \emph{farq}, \textquotedblleft{}a separation, a difference\textquotedblright{}, which reached Hindi through Persian.
\end{cousinweb}

\begin{grammarlens}[title={What you've built}]
One more word for saying what happens if.
\end{grammarlens}

\subsection*{Your turn}
\begin{itemize}
\raggedright
  \item \emph{Say it:} \emph{farq}
  \item \emph{Say it:} \emph{farq}, once more
  \item \emph{Say it:} \emph{agar tum roz ek ghaṇṭā paṛhoge to baṛā farq paṛegā.} — If you study for an hour every day, it will make a big difference.
  \item \emph{Recall:} say the Hindi for a magazine, then the Hindi for a scholarship, then say \emph{farq} again
\end{itemize}

\begin{tcolorbox}[breakable,colback=teal!4,colframe=teal!35!black,title={Before you move on}]
What does \hi{फ़र्क़} mean? (\textquotedblleft{}A difference\textquotedblright{}.) Say it once more.
\end{tcolorbox}
\section[\texorpdfstring{\hi{प्राप्त} \hi{करना} (prāpt karnā)}{प्राप्त करना (prāpt karnā)}]{\hi{प्राप्त} \hi{करना} (prāpt karnā) — to obtain, to receive (formal)}

\label{lesson:HI-C397-praptkarna}



\begin{quote}
Before the new one: say the Hindi for a scholarship, then the Hindi for a difference.
\end{quote}

\subsection*{You'll want to know: \hi{प्राप्त} \hi{करना}}
\textbf{\hi{प्राप्त} \hi{करना}} — \emph{prāpt karnā} — \textquotedblleft{}to obtain, to receive (formal)\textquotedblright{}.

\begin{quote}
\hi{अगर} \hi{तुम} \hi{परीक्षा} \hi{में} \hi{अच्छे} \hi{अंक} \hi{प्राप्त} \hi{करोगे} \hi{तो} \hi{तुम्हें} \hi{इनाम} \hi{मिलेगा।}
\end{quote}

\begin{quote}
\emph{agar tum parīkṣā meṁ acche aṅk prāpt karoge to tumhẽ inām milegā.}
\end{quote}

\begin{quote}
If you get good marks in the exam, you'll get a prize.
\end{quote}

Past, with \hi{ने} (he or she …): \hi{उसने} \hi{अच्छे} \hi{अंक} \hi{प्राप्त} \hi{किए}, \emph{usne acche aṅk prāpt kiye}.

\begin{cousinweb}
\textbf{\hi{प्राप्त} \hi{करना}} (\emph{prāpt karnā}) joins \hi{प्राप्त} to \hi{करना}, \textquotedblleft{}to do\textquotedblright{}. \hi{प्राप्त} is a tatsama, borrowed from Sanskrit \emph{prāpta}, \textquotedblleft{}attained, obtained\textquotedblright{} (Monier-Williams).
\end{cousinweb}

\begin{grammarlens}[title={What you've built}]
One more word for saying what happens if.
\end{grammarlens}

\subsection*{Your turn}
\begin{itemize}
\raggedright
  \item \emph{Say it:} \emph{prāpt karnā}
  \item \emph{Say it:} \emph{prāpt karnā}, once more
  \item \emph{Say it:} \emph{agar tum parīkṣā meṁ acche aṅk prāpt karoge to tumhẽ inām milegā.} — If you get good marks in the exam, you'll get a prize.
  \item \emph{Recall:} say the Hindi for a scholarship, then the Hindi for a difference, then say \emph{prāpt karnā} again
\end{itemize}

\begin{tcolorbox}[breakable,colback=teal!4,colframe=teal!35!black,title={Before you move on}]
What does \hi{प्राप्त} \hi{करना} mean? (\textquotedblleft{}To obtain, to receive (formal)\textquotedblright{}.) Say it once more.
\end{tcolorbox}
\section[(dialogue)]{Looking back: from Effort to Choosing What to Study}

\label{lesson:HI-R397-review-391-397}



\begin{quote}
Say the Hindi for an effort and toil.
\end{quote}

\subsection*{Your turn}
Say these aloud:

\begin{itemize}
\raggedright
  \item an effort, toil, enthusiasm, without a break, to repeat
  \item preparation, to try out, a method, experience, to occur to
  \item success, an achievement, to attain, a promotion, a matter of honour
  \item a prize, a competition, a post, capable, to put in order
  \item approval, an arrangement, to collect up, to fail, in the least
  \item upbringing, an ideal, generous, a good deed, to learn by rote
  \item a subject, a magazine, a scholarship, a difference, to obtain
\end{itemize}

\begin{tcolorbox}[breakable,colback=teal!4,colframe=teal!35!black,title={Before you move on}]
An effort? (\textbf{\hi{प्रयास}}, \emph{prayās}.) Toil? (\textbf{\hi{परिश्रम}}, \emph{pariśram}.) Enthusiasm? (\textbf{\hi{जोश}}, \emph{joś}.) Without a break? (\textbf{\hi{निरंतर}}, \emph{nirantar}.) To repeat? (\textbf{\hi{दोहराना}}, \emph{dohrānā}.) Preparation? (\textbf{\hi{तैयारी}}, \emph{taiyārī}.) To try out? (\textbf{\hi{आज़माना}}, \emph{āzmānā}.) A method? (\textbf{\hi{तरीक़ा}}, \emph{tarīqā}.) Experience? (\textbf{\hi{अनुभव}}, \emph{anubhav}.) To occur to? (\textbf{\hi{सूझना}}, \emph{sūjhnā}.) Success? (\textbf{\hi{सफलता}}, \emph{saphaltā}.) An achievement? (\textbf{\hi{उपलब्धि}}, \emph{upalabdhi}.) To attain? (\textbf{\hi{हासिल} \hi{करना}}, \emph{hāsil karnā}.) A promotion? (\textbf{\hi{तरक़्क़ी}}, \emph{taraqqī}.) A matter of honour? (\textbf{\hi{गौरव}}, \emph{gaurav}.) A prize? (\textbf{\hi{इनाम}}, \emph{inām}.) A competition? (\textbf{\hi{प्रतियोगिता}}, \emph{pratiyogitā}.) A post? (\textbf{\hi{पद}}, \emph{pad}.) Capable? (\textbf{\hi{क़ाबिल}}, \emph{qābil}.) To put in order? (\textbf{\hi{सँवारना}}, \emph{sãvārnā}.) Approval? (\textbf{\hi{मंज़ूरी}}, \emph{maṁzūrī}.) An arrangement? (\textbf{\hi{इंतज़ाम}}, \emph{intazām}.) To collect up? (\textbf{\hi{बटोरना}}, \emph{baṭornā}.) To fail? (\textbf{\hi{फ़ेल} \hi{होना}}, \emph{fel honā}.) In the least? (\textbf{\hi{क़तई}}, \emph{qataī}.) Upbringing? (\textbf{\hi{संस्कार}}, \emph{saṁskār}.) An ideal? (\textbf{\hi{आदर्श}}, \emph{ādarś}.) Generous? (\textbf{\hi{उदार}}, \emph{udār}.) A good deed? (\textbf{\hi{पुण्य}}, \emph{puṇya}.) To learn by rote? (\textbf{\hi{रटना}}, \emph{raṭnā}.) A subject? (\textbf{\hi{विषय}}, \emph{viṣay}.) A magazine? (\textbf{\hi{पत्रिका}}, \emph{patrikā}.) A scholarship? (\textbf{\hi{छात्रवृत्ति}}, \emph{chātravṛtti}.) A difference? (\textbf{\hi{फ़र्क़}}, \emph{farq}.) To obtain? (\textbf{\hi{प्राप्त} \hi{करना}}, \emph{prāpt karnā}.)
\end{tcolorbox}
\section[(dialogue)]{Choosing What to Study — once more}

\label{lesson:HI-R397-review-final-397}



\begin{quote}
Say the Hindi for a subject and a magazine.
\end{quote}

\subsection*{Your turn}
\begin{itemize}
\raggedright
  \item \emph{Say it:} a subject, a magazine, a scholarship, a difference, to obtain
\end{itemize}

\begin{tcolorbox}[breakable,colback=teal!4,colframe=teal!35!black,title={Before you move on}]
A subject? (\textbf{\hi{विषय}}, \emph{viṣay}.) A magazine? (\textbf{\hi{पत्रिका}}, \emph{patrikā}.) A scholarship? (\textbf{\hi{छात्रवृत्ति}}, \emph{chātravṛtti}.) A difference? (\textbf{\hi{फ़र्क़}}, \emph{farq}.) To obtain? (\textbf{\hi{प्राप्त} \hi{करना}}, \emph{prāpt karnā}.)
\end{tcolorbox}
